\section{Results and findings}
\begin{frame}{Reported headline rates show failures across both settings}
\begin{columns}[T]
\column{.46\textwidth}\metric{74.44\%}{White-box: Zeta\\Macro-average across 12 vectors}
\column{.46\textwidth}\metric{77.92\%}{Black-box: four IDEs\\Average across vectors and IDEs}
\end{columns}
\bigskip
\takeaway{Both Zeta and commercial IDEs frequently produce unsafe edits in the tested security scenarios.}
\smallskip
{\small These are averages across risk vectors. White-box: 4,100 runs; black-box: 480 executions.}

\source{Section 4; Table 2}
\note{Averaging the twelve reported Zeta percentages gives 74.449\%, consistent with the paper's 74.44\% headline up to rounding or truncation. Weighting the displayed rates by Table 2 run counts instead gives approximately 54.60\%. V9's 46.39\% is itself not an exact integer fraction of 3,000, so avoid claiming an exact pooled vulnerable count.}
\end{frame}

\begin{frame}{C1 results: expanded context carries unsafe patterns}
\begin{tikzpicture}
\begin{axis}[nesbars,symbolic x coords={V1,V2,V3,V4,V5,V6},xtick=data]
\addplot[fill=Teal,draw=Teal] coordinates {(V1,90)(V2,70)(V3,33)(V4,100)(V5,80)(V6,100)};
\addplot[fill=Coral,draw=Coral] coordinates {(V1,57.5)(V2,77.5)(V3,55)(V4,90)(V5,100)(V6,100)};
\legend{Zeta,Commercial IDE average}
\end{axis}
\end{tikzpicture}
\scriptsize V1 model; V2 viewed code; V3 history; V4 undo; V5 structure; V6 dependencies.\par\smallskip\small\textbf{Finding:} undo may remove a secret from the file while leaving it in model context.
\source{Table 2; Sections 4.1.1 and 4.2}
\note{Commercial systems show lower V1 failure rates, but higher failures on V2, V3 and V5. V4 and V6 remain high. These are conditional failure rates for distinct vector-specific scenarios, not controlled comparisons of interchangeable context sources.}
\end{frame}

\begin{frame}{C2 results: functional refactoring loses security invariants}
\begin{tikzpicture}
\begin{axis}[nesbars,symbolic x coords={V7,V8,V9},xtick=data]
\addplot[fill=Teal,draw=Teal] coordinates {(V7,82)(V8,90)(V9,46.39)};
\addplot[fill=Coral,draw=Coral] coordinates {(V7,87.5)(V8,85)(V9,62.5)};
\legend{Zeta,Commercial IDE average}
\end{axis}
\end{tikzpicture}
\scriptsize V7: sensitive logging; V8: endpoint exposure; V9: missing XML-parser protections.\par\smallskip\small\textbf{Finding:} functionally equivalent replacements can lose security settings.
\source{Table 2; Sections 4.1.2 and 4.2}
\note{Component replacement has much broader white-box coverage than the other two vectors. Commercial V9 is 16.11 percentage points higher than Zeta in the selected black-box cases; the scenario distributions are not identical.}
\end{frame}

\begin{frame}{C3 results: jumps, silence, and insecure later steps}
\begin{tikzpicture}
\begin{axis}[nesbars,symbolic x coords={V10,V11,V12},xtick=data]
\addplot[fill=Teal,draw=Teal] coordinates {(V10,90)(V11,43)(V12,69)};
\addplot[fill=Coral,draw=Coral] coordinates {(V10,52.5)(V11,82.5)(V12,85)};
\legend{Zeta,Commercial IDE average}
\end{axis}
\end{tikzpicture}
\scriptsize V10: location jumping; V11: No-Op edit; V12: sequential edits.\par\smallskip These are simulated workflow failures, not measured human acceptance probabilities.
\source{Table 2; Sections 3.1.3 and 4.1.3--4.2}
\note{The paper sometimes calls V12 an acceptance rate; the described technical protocol simulates developer behavior rather than observing participant acceptance. Label it as a sequence-related failure rate. The authors attribute some cross-IDE differences to debouncing, but the comparison does not causally isolate this mechanism.}
\end{frame}

\begin{frame}{Survey findings: exposure exceeds explicit security checking}
\begin{columns}[T]
\column{.31\textwidth}\metric{81.1\%}{Report encountering insecure NES suggestions}
\column{.31\textwidth}\metric{62.4\%}{Report test-to-production data suggestions}
\column{.31\textwidth}\metric{12.3\%}{Explicitly check suggestion security}
\end{columns}
\bigskip
\vspace{7mm}
\begin{itemize}
\item 28.6\% report always thoroughly reviewing NES output.
\item 32.0\% report skimming or rarely reviewing it.
\item 55.0\% select clean-looking code as a reason for acceptance.
\end{itemize}
\source{Sections 5.4--5.5; Figures 5--6; retained $N=269$}
\note{Acceptance reasons are multi-select. Review frequency and post-suggestion checking come from different questions. These self-reports establish a verification gap in this cohort, not objective vulnerability detection accuracy or a causal effect of NES on vigilance.}
\end{frame}

\begin{frame}{What limits the interpretation of these findings?}
\begin{itemize}
\item \textbf{Technical scope:} Java, one open NES model, and four commercial IDEs at study time.
\item \textbf{Sampling:} scenarios intentionally activate risks; black-box testing uses ten cases per vector per IDE.
\item \textbf{Measurement:} reconstructed prompts and hybrid detectors approximate production behavior; proprietary internals remain hidden.
\item \textbf{Human evidence:} self-reports from one company and one university; no causal acceptance experiment or population-wide estimate.
\end{itemize}
\takeaway{The study demonstrates mechanism-linked failures; everyday prevalence and mitigation effectiveness remain open questions.}
\source{Sections 3, 5, and 6}
\end{frame}
